\section{Related works}
\label{sec:related}

Mixture-of-Experts (MoE) architectures have emerged as a dominant approach for scaling LLMs beyond 100B parameters~\cite{deepseekv3, kimik22025}. By activating only a sparse subset of expert layers per token, MoEs decouple model capacity from compute cost, enabling efficient training at trillion-parameter scales.

However, this efficiency comes at the cost of increased memory footprint and expert-parallel communication overhead. To further reduce training cost, recent efforts have turned to low-precision computation~\cite{peng2023fp8}. Formats such as FP8 and FP4 not only improve compute throughput—NVIDIA GB200 NVL72, for example, supports up to 1440 PFLOPS of FP4 tensor operations compared to 720 PFLOPS for FP8 and 360 PFLOPS for BF16~\cite{gb200_specs}—but also reduce activation memory, enabling larger batch sizes and model scales within the same hardware footprint. When combined with systems-level communication frameworks like DeepEP~\cite{deepep2025}, which support expert-parallel data transfer in FP8 formats, low-precision has become a standard technique for efficient training of large language models.

DeepSeek-V3~\cite{deepseekv3} demonstrated the first successful deployment of blockwise-scaled FP8 training at the 671B MoE scale on Hopper GPUs, using its custom FP8 kernel library DeepGEMM~\cite{deepgemm2025} to achieve high arithmetic throughput. This work validated FP8 as a viable alternative to BF16 in production-scale training, delivering substantial improvements in compute efficiency and throughput. In parallel, NVIDIA introduced MXFP8, a microscaled FP8 format with native tensor core support on Blackwell GPUs. Training recipes based on MXFP8~\cite{mxfp8recipe} combine power-of-two block-level scaling with stochastic rounding to match BF16 convergence while significantly reducing activation and weight footprint. Subsequent work~\cite{fp8flowmoe2025} shows that excessive precision casts in BF16-dominated pipelines can erode FP8's expected gains, motivating careful end-to-end dataflow design.

To further improve training efficiency, recent work has explored 4-bit floating-point formats such as MXFP4 and NVFP4. MXFP4 is a standardized narrow-precision floating-point format defined in~\cite{ocp_mx_specification}, which introduces low-bit formats such as MXFP8, MXFP6, and MXFP4 for efficient AI training and inference. Tseng et al.~\cite{mxfp4training2025} demonstrated that MXFP4 can be used effectively in the backward pass by combining stochastic rounding (SR) with randomized Hadamard transforms (RHT), achieving near-lossless convergence on 6.7B-scale models. Meanwhile, NVFP4 introduces smaller block sizes and E4M3-style scaling, and has been used to train 12B models with near-FP8 accuracy using two-level quantization and format-aware rounding heuristics~\cite{nvfp4training2025}. Chmiel et al.~\cite{chmiel2025fp4wayfullyquantized} further demonstrated that fully quantized end-to-end FP4 training of LLMs is feasible with carefully designed scaling and optimization strategies.

These advances suggest that FP4 formats can substantially improve training efficiency, but they typically rely on hardware with native FP4 tensor core support and have not yet been validated at larger MoE scales. While some large-scale models such as GPT-OSS-120B~\cite{gptoss2025} are rumored to employ FP4 training, no public documentation is available regarding their training configurations.




\paragraph{Motivation}
Unlocking FP4-level efficiency without native hardware support would immediately benefit ongoing workloads—particularly Mixture-of-Experts (MoE) models, where expert-parallelism amplifies both activation memory and all-to-all communication costs. If software-based MXFP4 quantization could be integrated into existing Hopper-compatible pipelines without degrading accuracy or throughput, the potential payoff would be considerable: reduced memory footprint for larger batch sizes, lighter communication payloads across GPUs, and more aggressive activation checkpointing.

However, realizing this in practice introduces multiple technical challenges. Hopper's FP8-centric transformer kernels expect specific layout and scale formats, making naive insertion of FP4 quantization difficult. Intermediate data type conversions (e.g., FP8 $\leftrightarrow$ BF16 $\leftrightarrow$ FP4) incur latency, precision loss, and memory overhead. Prior work on FP8-based MoE training~\cite{fp8flowmoe2025} highlights that excessive precision casts can significantly erode the efficiency gains of low-precision operators, underscoring the need for careful dataflow and operator design when integrating lower-precision formats into existing pipelines. Additionally, MXFP4 employs per-block power-of-two scaling and E2M1 encoding, which are structurally incompatible with Hopper's FP8 E4M3 layout, requiring careful bit-level remapping and scaling alignment. The recomputation strategy must also be aware of asymmetric precision across forward and backward passes to avoid instability.

This work is driven by the need to address these system-level obstacles and fully realize MXFP4's potential on existing Hopper infrastructure. Our approach introduces a hybrid precision framework that compresses activations and all-to-all communication into MXFP4 format, while retaining FP8 for compute-intensive paths. By decoupling storage and arithmetic precision, and co-designing the quantization, dispatch, and dequantization flows, we achieve substantial efficiency gains without relying on hardware changes.